\section{有限 Markov 核与全变差工具}
\label{app:markov-tv}

本附录固定全文的行向量约定，并汇集后续证明所需的全变差估计。除非另有说明，本附录中的集合均非空且有限。

\subsection{核、确定性通道与观测算子}

从 $S$ 到 $T$ 的 Markov 核是函数
\[
  K:S\times T\longrightarrow[0,1],
  \qquad
  \sum_{t\in T}K(s,t)=1\quad(s\in S).
\]
全体此类核记为 $\Stoch(S,T)$。若 $K\in\Stoch(S,T)$、
$L\in\Stoch(T,U)$，则按
\[
  (KL)(s,u):=\sum_{t\in T}K(s,t)L(t,u)
\]
复合。概率分布视为单行 Markov 核，因此 $(\mu K)(t)=
\sum_s\mu(s)K(s,t)$。

映射 $q:S\to T$ 诱导确定性核
\[
  Q_q(s,t):=\one_{\{q(s)=t\}}.
\]
相应推前分布记为
\[
  q_*\mu:=\mu Q_q,
  \qquad
  (q_*\mu)(t)=\sum_{s:q(s)=t}\mu(s).
\]
若 $q$ 满射，则 $Q_q$ 的行恰好包含 $T$ 上的全部 Dirac 分布。对函数
$f:T\to\mathbb R$，核的观测算子定义为
\[
  (P_Kf)(s):=\sum_{t\in T}K(s,t)f(t).
\]
于是 $P_{KL}=P_KP_L$，并且 $P_{Q_q}f=f\circ q$。有限空间中，
$K$ 由 $P_K$ 唯一决定。

对 $A\subseteq T$，简记 $K(s,A):=\sum_{t\in A}K(s,t)$。确定性
粗粒化后的第 $s$ 行分布满足
\[
  (KQ_q)(s,\bar t)=K\bigl(s,q^{-1}(\bar t)\bigr).
\]

\subsection{全变差距离与行一致距离}

对 $T$ 上两个概率分布 $\mu,\nu$，定义
\begin{equation}
  \|\mu-\nu\|_{\TV}
  :=\frac12\sum_{t\in T}|\mu(t)-\nu(t)|.
  \label{eq:app-tv-def}
\end{equation}
有限空间上的等价表达为
\begin{equation}
  \|\mu-\nu\|_{\TV}
  =\max_{A\subseteq T}|\mu(A)-\nu(A)|.
  \label{eq:app-tv-events}
\end{equation}
特别地，全变差距离属于 $[0,1]$。对同型核 $A,B\in\Stoch(S,T)$，定义
\begin{equation}
  d_\infty(A,B)
  :=\max_{s\in S}\|A(s,\cdot)-B(s,\cdot)\|_{\TV}.
  \label{eq:app-dinfty}
\end{equation}
这是有限维核空间上的度量。它直接控制任意初始分布：
\begin{equation}
  \|\mu A-\mu B\|_{\TV}\le d_\infty(A,B).
  \label{eq:app-mixture-bound}
\end{equation}
本文以 $\Delta(T)$ 记 $T$ 上概率分布的单纯形。

\begin{lemma}[满射确定性前复合保持行最大值]
\label{lem:app-surjective-pullback}
若 $q:S\twoheadrightarrow T$ 且 $A,B\in\Stoch(T,U)$，则
\[
  d_\infty(Q_qA,Q_qB)=d_\infty(A,B).
\]
\end{lemma}

\begin{proof}
$Q_qA$ 在 $s$ 处的行就是 $A(q(s),\cdot)$；满射性保证对 $s$ 取最大值
与对 $t\in T$ 取最大值相同。
\end{proof}

\subsection{Dobrushin 系数}

下述系数源于 Dobrushin 对非齐次 Markov 链的工作；其作为全变差收缩系数的现代算子表述可参见\citet{Dobrushin1956,GaubertQu2015}。本文仍给出有限情形的自足证明，以固定行随机约定。

对 $R\in\Stoch(T,U)$，定义
\begin{equation}
  \alpha(R):=\max_{t,t'\in T}
  \|R(t,\cdot)-R(t',\cdot)\|_{\TV}\in[0,1].
  \label{eq:app-dobrushin}
\end{equation}

\begin{lemma}[全变差非扩张与 Dobrushin 收缩]
\label{lem:app-tv-contraction}
只要类型匹配，就有
\begin{align}
  d_\infty(CA,CB)&\le d_\infty(A,B),
  \label{eq:app-left-contraction}\\
  d_\infty(AR,BR)&\le \alpha(R)d_\infty(A,B),
  \label{eq:app-right-contraction}\\
  \alpha(RS)&\le\alpha(R)\alpha(S).
  \label{eq:app-alpha-submultiplicative}
\end{align}
\end{lemma}

\begin{proof}
对每个源状态 $x$，$(CA)(x,\cdot)$ 与 $(CB)(x,\cdot)$ 是分别以同一组
权重 $C(x,c)$ 混合 $A(c,\cdot)$ 与 $B(c,\cdot)$ 所得。由 $\ell^1$
范数的凸性，
\[
 \|(CA)(x,\cdot)-(CB)(x,\cdot)\|_{\TV}
 \le \sum_c C(x,c)\|A(c,\cdot)-B(c,\cdot)\|_{\TV},
\]
取最大值得到式~\eqref{eq:app-left-contraction}。

为证第二式，先取 $T$ 上概率分布 $\mu,\nu$，令
$r=\|\mu-\nu\|_{\TV}$。$r=0$ 时结论显然；否则令
\[
  \mu_+:=\frac{(\mu-\nu)_+}{r},
  \qquad
  \mu_-:=\frac{(\mu-\nu)_-}{r}.
\]
正负部分的总质量均为 $r$，故 $\mu_+,\mu_-$ 是概率分布，且
\[
 \mu R-\nu R
 =r\sum_{t,t'}\mu_+(t)\mu_-(t')
   \bigl(R(t,\cdot)-R(t',\cdot)\bigr).
\]
再用凸性与式~\eqref{eq:app-dobrushin}，得到
\[
  \|\mu R-\nu R\|_{\TV}
  \le\alpha(R)\|\mu-\nu\|_{\TV}.
\]
逐行取 $\mu=A(x,\cdot)$、$\nu=B(x,\cdot)$ 即得
式~\eqref{eq:app-right-contraction}。最后，将同一分布收缩不等式用于
$R(t,\cdot)$、$R(t',\cdot)$ 后取 $t,t'$ 的最大值，得到
式~\eqref{eq:app-alpha-submultiplicative}。
\end{proof}

\begin{remark}
若 $q:S\to T$ 是确定性映射，则 $\alpha(Q_q)=0$ 当且仅当 $q$ 为常值；
否则 $\alpha(Q_q)=1$。因此，确定性粗粒化虽然不会放大全变差误差，却通常
不会在这个最坏状态度量下产生严格收缩。
\end{remark}

后文所有近似结论都使用同一个 $d_\infty$。更换为平均分布加权距离会改变
量词和可见的坏状态；除非另行给定采样分布，二者不能互相替代。
